O executor experimental cria um projeto Docker Compose separado, sem utilizar as portas ou os volumes da aplicação interativa. Cada ensaio registra sua configuração, seus UUIDs e sua classificação. Os testes curtos de validação não entram na contagem de repetições definitivas. Os cenários planejados e os critérios prévios permanecem descritos no método; o catálogo de execuções distinguirá o que foi efetivamente realizado, as tentativas inválidas e seus motivos.

Na validação com PostgreSQL interrompido, foi programada carga de 5 eventos por segundo, com 5 segundos de aquecimento, intervalo de até 12 segundos para concluir o aquecimento, 18 segundos de medição e 5 segundos de interrupção solicitada ao controlador. O banco foi interrompido aproximadamente no primeiro terço da medição. Os instantes dos comandos de parada, da confirmação de parada e do comando de retomada foram registrados separadamente, pois a duração operacional pode diferir da pausa programada. Após a carga, foi aplicada uma observação de 5 segundos, além do tempo necessário à exportação.

A execução síncrona observou 91 eventos na fase medida: 39 respostas de sucesso e 52 respostas 503; o banco continha 39 desses registros ao final. A assíncrona observou 90 eventos: todos receberam 202, mas somente 89 estavam no banco, e um foi encontrado na DLQ por esgotamento das tentativas de persistência. Não havia aceite não conciliado nessas coletas. A mensagem na DLQ não desapareceu, mas também não equivale a um registro salvo. A diferença de uma iteração na fronteira das execuções é explicitada e impede tratar os totais como volumes exatamente iguais.

O comportamento observado é compatível com a distinção entre confirmação de publicação e conclusão da persistência. Ele impede afirmar que o aceite assíncrono garante que todo registro será salvo automaticamente depois que o banco voltar. Trata-se, contudo, de uma única validação curta por variante, não de uma estimativa de taxa de erro, capacidade ou recuperação generalizável. Não se confirmam ou refutam hipóteses comparativas de desempenho com esse ensaio.

Os perfis de rajada também foram exercitados como validação: a taxa passou de 5 para 10 eventos por segundo e retornou a 5, em uma janela de 18 segundos. Foram observados 120 registros na síncrona e 119 na assíncrona; todos esses registros foram conciliados com o banco, sem divergência de conteúdo ou hash. Os totais diferem nas fronteiras do agendamento e não representam uma comparação estatística de capacidade.

Na primeira validação com somente o consumidor interrompido, 91 eventos receberam 202, mas a observação inicial ainda encontrou 54 pendentes retidos e 37 salvos. Não se declarou perda. Um novo ensaio do mesmo perfil, com observação posterior de 40 segundos, conciliou todos os 91 eventos no banco, sem DLQ ou pendência final. Nesse segundo ensaio, 51 registros aceitos estavam sem confirmação de persistência no instante de conclusão do comando de retomada. O primeiro commit posterior foi observado aproximadamente 9,83 segundos após esse comando, e a persistência de todo esse conjunto pendente foi confirmada aproximadamente 10,44 segundos depois. Esses tempos não medem estabilidade sustentada de throughput e não são médias entre repetições. O aumento da janela de observação pertence à validação prévia dos instrumentos, não a uma alteração posterior de protocolo congelado.

A primeira tentativa de exportação do Kafka utilizou um consumidor de grupo cujo prazo de leitura terminou durante a atribuição inicial, produzindo exportação vazia apesar da confirmação da DLQ nos logs. A coleta foi rejeitada por inconsistência do instrumento e preservada. A exportação passou a atribuir diretamente a partição zero; um novo teste confirmou o conteúdo. Uma execução síncrona anterior também foi interrompida por alteração do script durante sua execução. Essas tentativas não foram incluídas como evidências válidas nem apagadas do histórico técnico.
